% NAL 2338, batch 2 — Gallica views 21–40.
% Pass: Opus 5.5 (claude-opus-5-5), 2026-09-23 — first pass, unchecked against the views by a human.
%
% The views. Greyscale scans (some stored as sRGB, all effectively grey),
% portrait, about 4316 × 5740–5756, one page per view; each leaf lies on a
% larger ground or on the next leaf, so the edge of another paper, and at
% times a few words of it, show at a top or side edge (v. 21, 22, 23, 25,
% 38, 40): not transcribed, noted. All twenty views carry text; none is
% skipped. Rectos and versos alternate, the recto carrying the number.
%
% What the twenty views hold. Three pieces, all Latin, all on the cycloid,
% the solids and spaces running to infinity, tangents and quadratures of
% the higher parabolas; all addressed to a « Vir Clarissimus » who is named
% Roberval (« Roberuallio ») in two headings.
%
%   v. 21–23   End of a letter, in a posed, regular copying hand. v. 21
%              opens mid-sentence on « demonstraui. », continuing the
%              page at v. 20 (batch 1). Viviani, Galileo, Cavalieri
%              (« Caualerio nostro ») and the weighing of cycloidal figures
%              cut from material; the « solidum hyperbolicum » of the
%              addressee. Ends « Flor. Kal. Octobr. 1643. » (v. 23), no
%              signature, no name of writer or addressee on these pages.
%   v. 24–35   A letter headed « Clar.mo Viro Roberuallio / Torr. S.P.D. »,
%              in a fast, heavy, sloping cursive with much ink bleed:
%              the priority of the trochoid, the centre of gravity of the
%              cycloid (a quotation of « Clar.mi Mersenni », v. 26), the
%              book de Motu, Desargues (« Librum de Sargues »), Guldin,
%              Cavalieri; then « Lemma I » (tangents of the higher
%              parabolas, v. 30–31), « Lemma II » (v. 31), two quadratures
%              of all parabolas (v. 32–34), the geometrical spiral (v. 34–
%              35), the hyperbola theorems sent to « Ill.m De Fermat ».
%              Ends « Vale. Dat. Florentię die 7 Julij ann. 1646 » (v. 35),
%              no signature, with a four-line postscript in paler ink.
%              Figures on v. 30, 31, 32, 33, described in notes.
%   v. 36–39   A separate piece, in the posed hand again: « Cl.mo Viro
%              Roberuallio / Lemma Pr.um », « Lemma secundum »,
%              « Theorema » (a figure « qualis a Cl. Roberual definita
%              est », by exhaustion), then the « Spiralis Geometrica »
%              (v. 37–38) and Michelangelo Ricci's spirals (« acutissimus
%              juuenis Michael Angelus Ricci », v. 39). Ends v. 39,
%              « demonstramus. », no date, no signature. At the top left of
%              v. 36, in another, cursive hand: « Hæc est demonstratio, quam
%              Torricellius in pag. 9.a et pag. 10 epistolæ ad Roberuallium
%              datæ 7.o die Julii, 1646. dicit à se missam fuisse in folio
%              separato. » — the only place in the batch where the name
%              Torricelli is written out. Its « pag. 9 », « pag. 10 » are the
%              bold numbers of v. 32–33, where the letter of v. 24–35 twice
%              promises a demonstration « in folio separato ».
%   v. 40      A second copy, in the posed hand, of the opening of the
%              letter of v. 24 (same heading « Torr. S.P.D. », same text,
%              a few spellings apart), with « 7. Julii 1646 » pencilled
%              (or in very pale ink) at the top left by another hand. It
%              runs off the foot of the page mid-sentence, « (quod satis
%              erat ad », into v. 41 (batch 3).
%
% Continuity. v. 21 continues v. 20 (batch 1's draft ends « Illam deinde
% quinquies diuersis semper principijs »; v. 21 begins « demonstraui. »).
% v. 40 runs on into v. 41. Inside the batch, v. 21→22 (the facing page's
% line ends show at the left of v. 22), 22→23, 24→25→…→35 (catchwords
% « tam », « coactus », « sum », « Iam », « DF » at the foot of versos),
% 36→37→38→39 run on.
%
% Who writes. The pages name no writer but « Torr. » in the salutation of
% v. 24 and v. 40, and the note of another hand on v. 36 names
% « Torricellius ». « Torr. » is transcribed as written and not expanded.
% The hands: (1) a posed, regular copying hand, v. 21–23 and 36–40 (the
% two are brought together by the letter forms, not proved one); (2) a
% fast, heavy cursive, v. 24–35, which writes the E of the figures as a
% large looped C (v. 30–33) and the 7 as an open chevron (v. 26, 35);
% (3) the reader's note at the top of v. 36; (4) the pale « 7. Julii 1646 »
% of v. 40. Nothing on these leaves is in a hand attributed to Mersenne,
% who is only named in them (v. 25, 26, 34, 40).
%
% Spelling and signs. The leaves' u/v, j/i, æ/ę/e are kept: hand (1) writes
% the æ ligature, hand (2) mostly e caudata (ę) and e. A long s is
% transcribed as s. Abbreviations kept: q; (-que), ꝑ (per, barred p),
% -b\textsuperscript{9} for -bus (whatever the exact shape of the sign),
% Clar.\textsuperscript{mo}, Cl., V. Cl.\textsuperscript{me}, gr. (grauitatis),
% cuboquad.\textsuperscript{m}, p\textsuperscript{m} (primum), \&, \&c.
% « n\textsuperscript{o} » renders an n with a raised loop read as « non »
% (v. 22, 28, 37, 39); batch 1 renders the same sign « nl ». A superscript
% curl over a final vowel (certè, propriò) is rendered as an accent.
% Letters of the figures are set in the text as written, not as maths. No
% notation of the period needs a LaTeX substitute.
%
% Numbers on the leaves.
% - A thin series at the top right of every recto, one per leaf, running
%   across all three pieces and all hands: 11 (v. 22), 12 (v. 24), 13
%   (v. 26), 14 (v. 28), 15 (v. 30), 16 (v. 32), 17 (v. 34), 18 (v. 36),
%   19 (v. 38), 20 (v. 40). No verso carries it (on v. 39 the 19 shows
%   through reversed). It behaves as the volume's foliation, but on this
%   microfilm it cannot be told pencil from ink, and no certificate is in
%   this batch; so no \folio{} is written, and the coordinator should set
%   it against batch 1's leaves and any certificate.
% - A bolder series beside it on the leaves of the letter of v. 24–35
%   only: a small cross « + » (v. 24, where the series puts 1), 3 (v. 26),
%   5 struck (v. 28), 7 (v. 30), 9 (v. 32), 10 at the top left of the
%   verso v. 33, 11 struck (v. 34). The pages of that letter; the note of
%   v. 36 cites its « pag. 9 » and « pag. 10 », which are v. 32 and 33.
%   Not a folio; notes only.
% - Dates on the leaves: « Flor. Kal. Octobr. 1643. » (v. 23); « Dat.
%   Florentię die 7 Julij ann. 1646 » (v. 35); « 7.o die Julii, 1646 » in
%   the note of v. 36; « 7. Julii 1646 » at the top of v. 40.
% - Other figures are content: « 7 ad 5 » (v. 26), « 5 », « 3 » (v. 30),
%   « 3 ad 2 » (v. 31), « 19 Quinti » (v. 32).
% - Stamps: round « BIBLIOTHEQUE NATIONALE / R.F. » (v. 24, 35, 36, 40).
%
% The catalogue gives no dating for this volume (empty in holdings.json),
% so \dating{} is left empty; no date is inferred here. The dates above
% are the leaves' own.

% Coordinator's reconciliation (2026-09-23), after all six batches:
%   One ink series of leaf numbers runs top right of the rectos across every
%   hand and piece: 1–10 (v. 2–20), 11–20 (v. 22–40), 21–29 (v. 42–59),
%   30–38 with « 34 bis » for the figure slip (v. 61–79), 39–48 (v. 81–99),
%   49–50 (v. 101–103). It ends on « 50 » at the last written leaf, as the
%   certificate of v. 1 has it (« Volume de 50 Feuillets / 1er Octobre
%   1888 »): it is the library's count. Being ink, it gets no \folio{}, in any
%   batch. The bold series on the letter of v. 24–35 (1–11) is that letter's
%   own pagination (batch 2).
%   Seams: v. 60 ends « vniuersa- » and v. 61 opens « lius »; batch 3 had set
%   the catchword into v. 60, and now ends where the page does. v. 100 runs
%   into v. 101 (« altius ; » / « in aliis autem humilius »), the « De Vacuo
%   Narratio » (v. 97–103). The « non » sign is set « nl » in batch 1 (hand B,
%   an n with a barred ascender) and « n^o » in batch 2 (the fast cursive, an
%   n with a raised loop): two hands, two signs, left as each pass read them.

\documentclass[11pt,a4paper]{article}
% The preamble shared by every transcription of the mathematicians' manuscripts in Gallica.
%
% It rests on one idea: **uncertainty compounds, it does not resolve.** A
% transcription that smooths over a doubtful word has destroyed the only thing
% it was made for — a reader must be able to tell, at every point, what was on
% the page and what was supplied. The apparatus macros below are therefore the
% critical apparatus, not decoration.
%
% The same commands are recognised by `scripts/render.mjs`, which produces the
% site's reading view, and by `scripts/tei.mjs`. A macro added here must be
% added there too, or rendering fails loudly — which is the wanted behaviour.
%
% Comments here are English, like the rest of the repository. The documents
% this preamble sets are French, and that is deliberate: see the skills.

\ProvidesPackage{germain}[2026/09/15 Transcriptions of mathematicians' manuscripts in Gallica]

\usepackage{amsmath,amssymb,amsthm}
\usepackage{mathrsfs}
% Kept from the parent preamble so the renderer's diagram support stays
% available; these manuscripts draw no commutative diagrams, and nothing here uses it.
\usepackage{tikz-cd}
\usepackage[english,french]{babel}
\usepackage{fontspec}

% --- Greek ------------------------------------------------------------------
% The text face stays fontspec's default, Latin Modern, which has no Greek:
% XeTeX drops every glyph it lacks with a line in the log and nothing on the
% page, and the Greek words the manuscripts quote vanished from the PDFs. So
% the Greek and Greek Extended blocks get a XeTeX character class of their own,
% and entering or leaving a run of it switches the family to CMU Serif — the
% same Computer Modern design, with polytonic Greek — and back. Only the
% family changes: a Greek word in \textit or \textbf keeps its shape and series.
% The transitions to and from every other class carry the tokens that class
% has towards an ordinary letter, so French spacing before « ; » or inside
% « » still holds after a Greek word. No grouping, so no brace or \endgroup
% can come out unbalanced; maths is left alone, since XeTeX inserts nothing
% there. Kept identical in `exercices.sty`.
\makeatletter
\newfontfamily\germainGreekfont{cmun}[
  Extension=.otf, NFSSFamily=germaingreek,
  UprightFont=*rm, ItalicFont=*ti, SlantedFont=*sl,
  BoldFont=*bx, BoldItalicFont=*bi, BoldSlantedFont=*bl]
\newXeTeXintercharclass\germain@greek
\def\germain@greekrange#1#2{%
  \@tempcnta=#1\relax
  \loop
    \XeTeXcharclass\@tempcnta=\germain@greek
    \ifnum\@tempcnta<#2\advance\@tempcnta\@ne\repeat}
\germain@greekrange{"0370}{"03FF}
\germain@greekrange{"1F00}{"1FFF}
\def\germain@greekprev{\rmdefault}
\def\germain@greekin{%
  \edef\germain@greekprev{\f@family}\fontfamily{germaingreek}\selectfont}
\def\germain@greekout{\fontfamily{\germain@greekprev}\selectfont}
\def\germain@greekjoin{%
  \XeTeXinterchartokenstate=\@ne
  \@tempcnta=\z@
  \loop
    \ifnum\@tempcnta=\germain@greek\else
      \XeTeXinterchartoks\germain@greek\@tempcnta=\expandafter{%
        \expandafter\germain@greekout\the\XeTeXinterchartoks\z@\@tempcnta}%
      \XeTeXinterchartoks\@tempcnta\germain@greek=\expandafter{%
        \the\XeTeXinterchartoks\@tempcnta\z@\germain@greekin}%
    \fi
    \ifnum\@tempcnta<4095\advance\@tempcnta\@ne\repeat}
% babel-french sets its own classes, and saves and restores the interchar
% state, each time a language is selected: join again after it does.
\addto\extrasfrench{\germain@greekjoin}
\addto\extrasenglish{\germain@greekjoin}
\AtBeginDocument{\germain@greekjoin}
\makeatother

\usepackage{geometry}
\geometry{a4paper,margin=25mm,marginparwidth=28mm}
\usepackage{marginnote}
\usepackage{xcolor}
\usepackage[normalem]{ulem}
\setlength{\emergencystretch}{3em}
\usepackage{hyperref}
\hypersetup{colorlinks=true,linkcolor=black,urlcolor=black,pdfborder={0 0 0}}

% --- Batch metadata ---------------------------------------------------------
% Taken as they stand by the rendering script, which shows them at the head of
% the reading view. They fix what the file claims to cover. The macro names are
% the parent project's, so that the scripts stay shared; \folder here means the
% volume's slug — `fr-9115` — and \pages counts Gallica views.

\newcommand{\folder}[1]{\gdef\grothFolder{#1}}
\newcommand{\batch}[1]{\gdef\grothBatch{#1}}
\newcommand{\pages}[2]{\gdef\grothFirst{#1}\gdef\grothLast{#2}}
\newcommand{\foldertitle}[1]{\gdef\grothFoldertitle{#1}}

% The holder's own shelfmark, printed as they write it: « Français 9115 ».
\newcommand{\shelfmark}[1]{\gdef\grothShelfmark{#1}}

% The dating, copied verbatim from the holder's catalogue — never inferred
% here. « XIXe siècle » is the BnF's; a pass that dates a leaf from its content
% says so in a \note{}, as its own inference.
\newcommand{\dating}[1]{\gdef\grothDating{#1}}

% The status notice, printed in the title block and repeated as a diagonal
% watermark on every page of the PDF. The manuscripts are in the public
% domain; what this line declares is not a right but a quality — an
% unverified first pass by a machine. The skills write the line; a file
% without it gets no watermark, so leaving it out is visible in the output.
\newcommand{\watermark}[1]{\gdef\grothWatermark{#1}}

\gdef\grothFolder{?}\gdef\grothBatch{}\gdef\grothFirst{?}\gdef\grothLast{?}%
\gdef\grothFoldertitle{}\gdef\grothShelfmark{}\gdef\grothDating{}\gdef\grothWatermark{}

% --- The résumé -------------------------------------------------------------
\newenvironment{resume}
  {\small\setlength{\parskip}{0.45em}}
  {\par\medskip\hrule\medskip}

% --- Keywords ---------------------------------------------------------------
% In English, closing the modernised reading's résumé: the modern vocabulary
% under which to search for what the leaves construct. The manifest extracts
% them to tag the volume — this line is the single source of the tags.
\newcommand{\keywords}[1]{\par\smallskip\noindent
  {\footnotesize\textsc{Keywords} --- \textit{#1}}\par}

% --- The critical apparatus -------------------------------------------------

% The Gallica view number, in the margin. This is the anchor that lets a
% disagreement over a formula be settled by looking at *one* image rather than
% a volume of 750. The site uses it to turn the facsimile as the transcription
% is scrolled. A view is one scanned image: usually one side of a leaf.
\newcommand{\page}[1]{%
  \marginnote{\footnotesize\color{gray!70}\textsf{v.\,#1}}%
  \label{page:#1}\ignorespaces}

% The BnF's pencilled foliation, where the leaf carries it and a pass can read
% it: « 198r ». This is what the literature cites — « ff. 198r–208v » — and
% the one line that turns a citation into a view. Recorded, never inferred:
% a folio a pass cannot read is not written.
\newcommand{\folio}[1]{%
  \marginnote{\footnotesize\color{gray!50}\textsf{f.\,#1}}\ignorespaces}

% The modernised reading's coarser counterpart to \page{}: which views a
% section is drawn from.
\newcommand{\pagerange}[2]{%
  \marginnote{\footnotesize\color{gray!70}\textsf{v.\,#1--#2}}%
  \ignorespaces}

% Illegible: never guessed.
\newcommand{\ill}{\textcolor{red!60!black}{[\,\ldots\,]}}

% A doubtful reading: the word is given, and flagged as uncertain.
\newcommand{\uncertain}[1]{\uline{#1}}

% An editorial addition — a missing letter, an implied word.
\newcommand{\add}[1]{\textcolor{blue!50!black}{[#1]}}

% Struck out by the writer. What was crossed out is part of the document.
\newcommand{\struck}[1]{\sout{#1}}

% The transcriber's note, on the state of the page or the mathematics.
\newcommand{\note}[1]{\par\smallskip{\small\color{gray!60!black}\textit{[#1]}}\par\smallskip}

% A marginal note of hers, distinct from the body of the text.
\newcommand{\marginal}[1]{\par\smallskip\noindent\hspace{1em}%
  {\small\color{gray!60!black}#1}\par\smallskip}

% --- Title ------------------------------------------------------------------
% The title block is French, like the documents themselves.

\AddToHook{shipout/background}{%
  \ifx\grothWatermark\empty\else
    \begin{tikzpicture}[remember picture, overlay]
      \node[rotate=52, scale=3.4, align=center, text=gray!18,
            font=\sffamily\bfseries]
        at (current page.center) {\grothWatermark};
    \end{tikzpicture}%
  \fi}

\AtBeginDocument{%
  \begin{center}
    {\large\bfseries Manuscrits de mathématiciens (Gallica) ---
      \ifx\grothShelfmark\empty\grothFolder\else\grothShelfmark\fi}\\[2pt]
    {\itshape\grothFoldertitle}\\[4pt]
    {\small vues \grothFirst--\grothLast
      \ifx\grothBatch\empty\else\ (lot \grothBatch)\fi}\\[2pt]
    \ifx\grothDating\empty\else
      {\small\color{gray!60!black}Datation du catalogue : \grothDating}\\[2pt]
    \fi
    \ifx\grothWatermark\empty\else
      {\small\bfseries\color{red!45!gray}\grothWatermark}\\[2pt]
    \fi
    {\footnotesize\color{gray!60!black}Transcription établie d'après les images IIIF de Gallica.
     Source gallica.bnf.fr / Bibliothèque nationale de France.\\
     Les lectures incertaines sont soulignées, les illisibles notées [\,\ldots\,],
     les ajouts entre crochets ; \textsf{v.} numérote les vues de Gallica,
     \textsf{f.} la foliotation au crayon de la BnF.}
  \end{center}
  \vspace{1em}}

\endinput


\folder{nal-2338}
\shelfmark{NAL 2338}
\batch{2}
\pages{21}{40}
\dating{}
\watermark{Lecture automatique\\non vérifiée}
\foldertitle{Papiers divers de Mersenne.}

\begin{document}

\page{21}
\note{Verso, sans numéro ; la page commence au milieu d'une phrase, qui continue la vue précédente (lot 1). Au-dessus du bord supérieur de la feuille dépasse la tranche d'un autre papier, portant quelques lettres d'une autre écriture, qui n'appartiennent pas à cette page ; la page du recto transparaît à l'envers. Main de copie posée, régulière, en latin, qui est celle de toute cette lettre (vues 21–23).}

demonstraui. Quo ad solida nihil habeo. Tangentem prædictæ lineæ iam ostenderat mihi Vincencius Viuianus florent: Clar.\textsuperscript{mi} Galilei alumnus, etiamnunc adolescens. Quo ad Auctorem hujus figuræ, credo ego ingenium suum accutissimum, et feracissimum illam ex se obseruare potuisse nemine indicante ; hujusmodi enim linea naturæ familiaris erat ; constatque ex compositione duorum motuum recte, et circularis ; Attamen viuunt adhuc testes, quibus olim Gal.\textsuperscript{s} irritas Lucubrationes suas comunicauit circa hanc figuram ; imò supersunt paginæ aliquot Clar.\textsuperscript{mi} Mathematici, in quibus et picturas, et \uncertain{aggrestiones} suas nonnullas circa hoc subjectum cùm adolescens delineauerat. Pluribus ab hinc annis Theorema hoc proposuit ille mirabili Geometræ Caualerio nostro, ipsique dixit idem quod et mihi, et pluribus alijs confirmauit, nempe se olim experimentum fecisse, appensis ad Libellam spatijs figurarum materialibus, quantuplum esset cycloidale spatium ad circulum suum genitorem ; et semper illud inuenisse, nescio quo facto, minus quàm triplum ; ideoque incæptam contemplationem deseruisse, ob incommensurabilitatis suspicionem. Quid si aliquandò, inconstanti fallacia, reperisset minus quàm triplum, aliquandò verò majus, tunc asserebat Lyncæus Mathemat.\textsuperscript{cus} vlteriorem contemplationem prosecuturũ fuisse ; rejectâ scilicet variationis causâ in materiæ inequalitate, atque raturę.
\note{« suas » (après « Lucubrationes ») et « paginæ » sont d'une encre plus pâle, comme remplis après coup dans un blanc laissé à la copie. « aggrestiones » : lu tel quel, avec un petit signe sur le e ; on peut y voir « aggressiones », si « st » est un s long suivi d'un s. « Lyncæus » : l'æ est lié. « raturę » : e à queue (e caudata). L'orthographe « accutissimum », « comunicauit » est celle de la page.}

Propositionem illam de solido a qualibet coni sectione circa axem reuoluta descripto, atque de ejusdem solidi
\note{La page s'arrête au milieu de la phrase, au tiers inférieur ; le reste est blanc. Les traits horizontaux en fin de ligne sont des traits de remplissage, non des ratures.}

\page{22}
\note{Recto. En haut à droite, d'un trait fin à l'encre, « 11 » (voir l'en-tête). À gauche de la vue paraissent les fins de ligne de la page de la vue 21, qui lui fait face : les deux pages se suivent, et la phrase interrompue à la vue 21 continue ici. Même main.}

centro grauitatis, vnica simul breuique demonstratione ostendimus ; suppositâ tantùm modica Apollonij cognitione. Verum duplex hoc Theorema inter neglecta a me rejicitur ; nullum enim habebit locum in opusculis, quæ nunc \uncertain{propolare} cogor, in quibus præcipuè profiteor materiæ vnitatem.
\note{« propolare » : lu tel quel ; la deuxième voyelle se lit o plutôt que a.}

Quò ad solidum hyperbolicum, iam n\textsuperscript{o} meum, sed tuum, Dispeream si iam amplius spero me visurum tam sublimem, et tam doctam demonstrationem, quæ cum tua conferri mereatur. Opptimum equidem maximumque nunc percipio laborum meorum fructum, eo tantum nominè, quòd tu vir Clar.\textsuperscript{me} atque ingeniosiss.\textsuperscript{me}, tam acutis demonstrationib\textsuperscript{9}. tantâque doctrinæ affluentia vnicam ineptiolam meam illustrare dignatus sit. Gratias primùm ago maximas ; Deinde vt desiderio tuo satisfaciam, methodus mea circa Demonstrationem hujus solidi, diuersissima est à tua. Altera quidem ex meis aggressionib\textsuperscript{9} ꝑ doctrinam indiuisibilium procedit, quæ si cum erudito lectore semper ageretur, paucissimis verbis expediri posset. Altera verò ꝑ inscriptionem, et circumscriptionem more veterum, non adeò expedita est, sed facilis, et fortasse curiosa. Hoc vnum reperi in tua Scriptura, quod conueniat cum meis, nempe constructio illa prosecando frusto solidi hyperbolici in data ratione. Demonstrationes verò ab eadem constructione dissimillimæ emanant.
\note{« n\textsuperscript{o} » : n suivi d'une petite boucle levée, abréviation de « non ». « demonstrationib\textsuperscript{9} », « aggressionib\textsuperscript{9} » : b suivi du signe de -us (voir l'en-tête). « ꝑ » : p barré sous la ligne, « per ». « Quò » : un accent sur l'o. « Opptimum » est l'orthographe de la page.}

Cæterum euidentiores agnosco hyperbolas in Laudibus quibus me exornas, quàm in demonstrationibus quibus hyperbolicũ

\page{23}
\note{Verso du feuillet numéroté « 11 », sans numéro. La phrase de la vue 22 continue ici. Au-dessus du bord de la feuille dépasse, comme à la vue 21, la tranche d'un papier plus grand, avec quelques lettres d'une autre écriture qui n'appartiennent pas à cette page. Même main.}

solidum ipse metiris. Vtinam illis aliquandò dignus fiam, vt in lectione operum tuorum, quæ auidissimus expecto, illa intelligere valeâm, fructusque scientiæ suauissimos, et diuitias ingenij inestimabiles inde colligere possim, et intellectum meum ditare. Vale Vir Clar.\textsuperscript{me} tuorumque operum editionem accelera in publicam Litteratorum omnium vtilitatem.

Flor. Kal. Octobr. 1643.
\note{La date, en grands caractères, termine la lettre ; « Flor. » abrège sans doute « Florentiæ » (inférence). Aucune signature, aucun nom de l'auteur ni du destinataire sur les pages de ce lot ; la lettre s'adresse à un « Vir Clar.\textsuperscript{me} » qui a démontré la mesure du « solidum hyperbolicum », et parle de Galilée, de Viviani et de Cavalieri (« Caualerio nostro »). Le reste de la page est blanc.}

\page{24}
\note{Recto d'une nouvelle feuille, qui commence une nouvelle lettre ; la page est écrite jusqu'au pied. En haut à droite, à l'encre, « 12 » suivi d'une petite croix (voir l'en-tête). Autre main que celle des vues 21–23 : cursive rapide, penchée, à longues hampes et à pleins très appuyés ; l'encre a beaucoup bavé, et des taches d'encre et des décharges de la page voisine couvrent la moitié droite du haut de la page. Cachet rond « BIBLIOTHEQUE NATIONALE » avec « R.F. » au milieu de la page.}

Clar.\textsuperscript{mo} Viro Roberuallio

Torr. S.P.D.
\note{La suscription se lit telle quelle : « Torr. » est l'abréviation d'un nom, écrite sur la feuille ; son développement n'est pas donné ici.}

De Trochoide (esto enim quantumlibet Trochoides) \uncertain{siue} Italicum siue Gallicum Problema sit, nihil mea interest. Meum certè non est quod ad inuentionem attinet. De autore ipsius quod ego acceperam ab amicis illud credideram scripseramq; Vos aliter vultis ? ꝑ me iam licet. Hoc certissimum est (quicquid tandem feratur ab alijs) celeberrimum Galileum usq; ad supremum vitę \uncertain{diem} mensuram illius figurę ignorauisse quam ex Gallia non accepit, ubi fortasse inuenta non fuerat : illud certè profitebatur ; neque uideo cur demonstrationem illius si à quopiam accepisset in commune non protulisset ad gloriam suam, quamquam aliena esset. Ego fateor non adeo multis ab hinc annis demonstrationes illas me reperisse, sed propriò Marte non minus quam à quopiam alio, siue ante me, siue post, factum \struck{est} sit. Si uero aliqua ex meis demonstrationib\textsuperscript{9} conuenit cum Gallicis, primùm quod ad meam internam quietem attinet, quodq; plurimi facio, ego mihi conscius sum illas omnes ex meo reperisse, \struck{\&} et quicunq; me nouerit idem credet : deinde quicquid alij credant, nihil me mouet. Eximiũ illum voluptatis fructum quem percipimus unusquisque in inuentione ueritatis, et pro quo tantum speculor, nemo à me auferet. De gloria quam ꝑ contentiones et controuersias acquirere debeam minimè sollicitus sum : Propterea non tantum unam, sed et omnes demonstrationes illas, si quis volet concedere paratus ero, dumodo ꝑ iniuriam non eripiat. Sed de Centro grauitatis Cycloidis
\note{« siue » (le premier) est sous une tache. « diem » s'écrit ici comme « deem » surmonté de points. « Vos aliter vultis ? » : le signe qui suit « vultis » est lu comme un point d'interrogation. « sit » est écrit au-dessus de « est », biffé. Dans « \struck{\&} et », un signe biffé avant « et ». Plusieurs mots portent un signe levé au-dessus de la dernière lettre (« certè », « propriò », « minimè ») : il est rendu par un accent. La phrase continue à la vue suivante.}

\page{25}
\note{Verso de la feuille « 12 », sans numéro ; même main rapide, mêmes taches d'encre en haut à gauche. La phrase de la vue 24 continue ici. Sous le bord inférieur de la feuille paraît le haut d'un autre papier, écrit d'une autre main, qui n'appartient pas à cette page.}

\uncertain{scis} profectò V. Cl.\textsuperscript{me} demonstrationem illius misisse in Galliam (atq; utinam non misissem) precibus Cl. Mersenni integro biennio antequam illud habere diceres, ut in ultimis tandem epistolis habuisse iam diu confiteris. In illa demonstratione mea ostendebatur à me dato Centro grauitatis et mensura alicuius plani (quod satis erat ad intentum meum) ipsius solidum demonstraui. Ipse vero dicis V. Cl. tandem habere Methodum pro reperiendo centro gr. plani ex data solidi planiq; mensura. Propositiones conuersę sunt. Inuersio autem huiusmodi facillima est ; et si in alterũ ex nobis suspicio aliqua ferri debeat, certè in me non cadet ; nam multo antequam de hoc uerbum faceres cum nostris Italis, sed et cum Gallis non solum enunciationem, sed etiam demonstrationem ipsam, orantibus uobis, ego misi in Galliam. Illud etiam pro me stare uidetur argumentum, quod numquam ne uerbum quidem fecisti de Centro gr. Cycloidis cum interea tantoperè, et quidem meritò, \uncertain{gloriareris} de omnibus alijs, Quadratura, Tangentibus, solidis etiam de eo circa axem, quod tantùm sperare dicebas. Veri simile non est cum reliqua omnia proponeres, et in lucem edere velle promitteres quadraturas, tangentes, et solida, de unico Centro grauitatis siluisse, si illud tantum sperauisses, quod quidem Problema meo iudicio nulli reliquorum posthabendum uidetur. Sed de his, si opus fuerit multo plura dicemus suo temporè. Nemo
\note{« scis » : le premier mot, sous une tache, se lit ainsi. « gloriareris » : le mot est écrit serré et se lit aussi « gloriareis ». Au pied, sous la dernière ligne, la réclame « tam », d'un trait plus fin : le mot qui commence la page suivante.}

\page{26}
\note{Recto d'une nouvelle feuille, de même format et de même main que la feuille « 12 ». En haut à droite, deux nombres : « 13 », d'un trait fin, et à sa droite un « 3 » plus grand et plus gras (voir l'en-tête). Au-dessus de la feuille paraît le papier plus grand sur lequel elle est posée, blanc. La réclame « tam » de la vue 25 est le premier mot de la page : la lettre continue.}

tam facilè suam laudem unicuiq; tribuet quam ego, dummodo tamen ignorantia, uel credulitate non decipiar. Methodum pro tangentib\textsuperscript{9} ex doctrina motus ego reperi pluribus ab hinc annis, nulla ab alijs habita luce, uel auxilio : Cum amicis contuli ; et in multis figuris propagaui. Postea incidi in demonstrationes Trochoidis et utrumq; uulgaui inter amicos, antequam in meis libellis ederem. Ex improuiso, quando nil tale sperabam, nuncius horribilis ex uobis affertur hęc omnia ante me uos etiam inuenisse. Si uerum hoc est, certè pro meis illa amplius non essent habenda (quamquam fortasse nullus mortalium ad hęc umquam descenderet). Vide Vir Clarissime quam ingenuè ego agam \uncertain{cedendo}, etiam ea quę iure ęque mea sunt, ac uestra, cum uterq; proprio Marte adinuenerit, abstracta, (si qua intercesserit) modici temporis differentia. Sed incredibile est quanta iniuria afficiar dum uideo mihi præripi ea, quę mea esse deberent sine controuersia. Inuoco hominum fidem. Ecce uerba ipsissima uestri Clar.\textsuperscript{mi} Mersenni in Epistola ad me data postquam enunciationem tantum Centri gr. Cycloidis sine demonstratione ad uos miseram (dicebam enim secare axem in ratione \uncertain{7} ad 5.)
\note{« cedendo » : le mot se lit aussi « credendo » ; la deuxième lettre n'est pas nette. « præripi » : æ lié. « 7 » : chiffre en forme de chevron ouvert à gauche, la forme italienne du 7 ; le sens n'est pas vérifié ici. La parenthèse n'est pas fermée sur la page ; la phrase s'arrête sur « 5. ».}

Dubitat Noster Roberuallus an Mechanicè tantum centra gr. Cycloidis et Semicycloidis inuenerit, quę Geometricè falsa suspicatur ; docebis num istius rei demonstrationem habeas.
\note{Les mots « quę Geometricè falsa suspicatur » sont soulignés par l'auteur de la lettre. Une tache d'encre couvre en partie « falsa ». Le paragraphe, introduit comme les « uerba ipsissima » de Mersenne, est la citation ; il n'y a pas de guillemets sur la page. Le reste de la page est blanc.}

\page{27}
\note{Verso de la feuille « 13 », sans numéro ; même main. La page de la vue 26 s'achevait sur la citation ; celle-ci reprend la parole de l'auteur de la lettre.}

Quare ergo V. Clar.\textsuperscript{me} dubitabas et Geometricè falsum suspicabaris, quod ipse sciebas. Non dubito ego an uestra plana in infinitum abeuntia figurę genitrici ęqualia sint, nam demonstrationem habeo. At ego non docui num demonstrationem illam haberem, sed protinus (quod quispiam præter me non fecisset) acceptam demonstrationem in Galliam misi, simulque methodum uerè meam, cum demonstratione pro inueniendo alterutro siue centro, siue solido in omnibus figuris, ex altero dato cum quadratura : quam methodum si apud uos antè habebatis, certè de ueritate illa geometrica non erat quod dubitaretis, nam solidum circa basim tunc temporis uos habuisse dicitis. Sed missa faciamus hęc : Habeo enim et aliam methodum, quę unica enunciatione determinat reperitq; Centrum grauitatis linearum, superficierum ex reuolutione natarum, planorum, corporumq; omnium, dummodo axem, siue diametrum habeant. Vulgata est hęc apud amicos Italos ; Oro uos ne inter uestra hanc etiam habeatis, nam hoc esset tollere penitus omne literarum scientiarumq; commercium. De libellis meis si tibi uir Clar.\textsuperscript{me} placuerunt ut ipse affirmas (excepto tamen de Motu et Trochoide) maximas ago gratias. In hoc autem plurimum discrepamus, nam mihi ex ipsis nullus placet. Meliora non edidi, quia meliora habere non credebam : talia edidi, quasi
\note{« abeuntia » : les lettres « nt » sont récrites en surcharge, à l'encre plus épaisse. « sciebas » est suivi d'un point surmonté d'un trait, lu comme un point. Au pied, la réclame « coactus », d'un trait fin.}

\page{28}
\note{Recto. En haut à droite, « 14 », d'un trait fin, et à sa droite un « 5 » plus gras, barré d'un trait oblique (voir l'en-tête). Même main. La réclame « coactus » de la vue 27 est le premier mot : la lettre continue.}

coactus. Videbar enim quodammodo ad hoc teneri ex titulo officij mei. Vtinam commercium hoc cum Clar. Galliarum Geometris accidisset mihi ante aliquos annos, nam certè uel meliora edidissem, uel nulla. Si reliqua bona essent leuis esset iactura libellum de Motu et quidem totum ex meis opusculis abolere. Argumenta quę producis V. Clar.\textsuperscript{me} contra fundamenta illius doctrinę fortia sunt, sed tamen non sunt noua. Extant huiusmodi obiectiones, immo etiam et multò fortiores apud ipsum Galileum in libro Proiectorum, sed latitant in illis Dialogis, quos ipse inter latinas Propositiones Italicè interserit. Respondere possem eadem, quę respondet Galileus, et non nihil etiam ex meo ; sed non ęstimo tanti illum meum libellum, ut hanc molestiam susciperem tecumq; eandem partiri V. Clar.\textsuperscript{me} ęquum censeo. Ad ultimum nemo hanc defensionem mihi eripiet, si dixero libellum esse ex maiori parte Geometricum. Archimedes supposuit olim proiecta non ꝑ parabolas, sed ꝑ lineas spirales suas procedere, et huius suppositionis gratiâ totum libellum composuit. At suppositionem esse falsam non multò post \uncertain{comperiit}. Quid faceret. librum \uncertain{n\textsuperscript{o}} illum damnare penitus debebat ? quin potius eadem omnia legi sineret sublata omni Proiectorum memoria, additoq; uocabulo puncti cuiusdam, quod non naturali sed ficta quadam motuum lege moueatur. Proprium habet hoc Geometria quod nullis alienis auxilijs suffulta,
\note{Sous « Italicè » et « inter », un petit signe en croix, d'une autre encre, peut-être une marque de lecteur. « interserit » se lit aussi « interseret ». « comperiit » : le mot se lit « compereet » surmonté d'un accent ; lecture proposée avec doute. « n\textsuperscript{o} » : n suivi d'une boucle levée, lu comme l'abréviation de « non ». « totum » (dans « gratiâ totum libellum ») est d'un trait plus fin. La phrase continue à la page suivante ; le bas de la page est blanc.}

\page{29}
\note{Verso de la feuille « 14 », sans numéro ; même main. La phrase de la vue 28 continue ici.}

sibi ipsi tantum innixa numquam ruere potest, sed mole sua stat, licet reliqua omnia, quę propria culpa ipsi malè adhęrebant, dilabantur. Reiciamus ergo quicquid Physicum est in eo libello, nempe uocabula proiectorum, grauium, ballistarum et.c. et Geometricum maneat, hoc est propositiones abstractę. Cęterum Tabulę, instrumenta, et si quid aliud huiusmodi est, non sint pro mensurandis proiectorum iactibus, sed pro determinandis quibusdam lineis geometricis, alijsq; ad geometricas parabolas pertinentibus. Librum de Sargues numquam uidi, sed quomodo probet ea quę proponit nescio. Videas ipse an liceat unicuiq; talia scribere, et fortasse etiam pari ratione probare. Videas autem ipse V. Clar.\textsuperscript{me} quam facilè sit hoc \uncertain{euo} Scriptorum feracissimo in aliorum reperta incidere, et pro suis innocenter aliena circumferre. Ipse proponis tamquam nouum quod Quadratrix transeat ꝑ Centrum grauitatis semiperipherię et.c. At multis ab hinc annis hoc ipsum demonstratum fuit, et editum ab alijs multò uniuersalius.
\note{« de Sargues » est écrit en deux mots. « euo » : lu tel quel (pour « æuo » ?), avec doute sur la première lettre.}

Hoc ipsum ostenderam ego quoque maxima, et breuitate, et simplicitate, antequam scirem quemdam Guldinum Iesuitam hęc edidisse. Monitus enim fui nuperrime a Clar.\textsuperscript{mo} Viro Caualerio me in idem inuentum incidisse atq; adeo perdidi contemplationem qua nullam habeo elegantiorem, aut magis absolutam. Admiratus
\note{Au pied, la réclame « sum ». Le bas de la feuille est blanc.}

\page{30}
\note{Recto. En haut à droite, « 15 », d'un trait fin, et à sa droite un « 7 » plus gras (voir l'en-tête). Même main. La réclame « sum » de la vue 29 est le premier mot : la lettre continue.}

sum subtilissimum Theorema de spatijs in infinitam longitudinem abeuntibus. Egregium inuentum, et quod \uncertain{magis} meis uel ipse preferam. At non dico inter mea illud habuisse cum ad infinitam illam distantiam (exceptis hyperbolis meis) non \uncertain{respexerim}. Circa illius demonstrationem nihil prorsus moratus sum ; nam uix lecta propositione statim animaduerti ipsi conuenire demonstrationem, qua iam pridem usus fueram in dimensione infinitarum parabolarum. Vulgata iam est apud amicos Italos, ipsamq; hic obiter adiungam orsus a Tangentibus, quandoquidem tangentes in ea adhibentur. quę omnia (licet ineptię sint) fidei uestrę commendo.
\note{« magis » est récrit en surcharge sur un autre mot. « respexerim » : le début du mot est surchargé et empâté ; on peut lire aussi « inspexerim ».}

Lemma I.

Esto parabola quęlibet ABC ad cuius punctum A danda sit tangens. Applicetur AD ad diametrum ED, fiatq; ut exponens applicatarum quę exemp: gr: sit 5 ad exponentem diametralium quę exemp: gr: sit 3 ita recta ED ad DC dico iunctam AE non conuenire cum parabola ad aliud punctum præter A. Nam si potest conueniat in B, et applicata BI fiat ut ED ad DC, ita EI ad IO sitq; latus rectum CL, cui ęquidistet OP.
\note{Figure dans la marge droite : un axe vertical, qui porte de haut en bas les points E (en minuscule, au sommet), O, C, I (en minuscule), D ; la droite EA descend à gauche jusqu'à A ; une seconde droite descend à droite de E par P et L ; la parabole part de C, où la droite CL est menée perpendiculairement à l'axe, et passe par B jusqu'à A ; les ordonnées OP, CL, BI et AD. La lettre E du texte est un grand C bouclé, distinct de C ; on la lit E d'après la figure. Le texte écrit « diametrum ED » : la seconde lettre est lue D, la première est ce E.}

Erit iam dignitas BI, hoc est in nostro casu cuboquad.\textsuperscript{m} BI ęquale ductui dignitatum IC in CL, nempe ductui Cubi IC in quadratum CL ob latus rectum. At idem cuboquad.\textsuperscript{m} BI ęquatur ductui cubi IO in quadratum OP (ut mox patebit) ergo ductus IC in CL, et IO in OP sunt ęquales. Quod est impossibile, ut demonstrabimus inferius.
\note{« cuboquad.\textsuperscript{m} » : « cuboquadratum », la cinquième puissance, écrit en abrégé. Le bas de la page est blanc.}

\page{31}
\note{Verso de la feuille « 15 », sans numéro ; même main. La démonstration du Lemma I continue.}

Quod promisimus primò ostendetur sic. Cuboquadratum AD ad cuboquad.\textsuperscript{m} BI rationem habet compositam ex ratione cubi AD, ad BI, siue cubi DE ad EI, siue cubi DC ad IO. Et ex ratione quadrati AD ad BI siue quad.\textsuperscript{ti} DE ad EI uel quadrati CE ad EO, hoc est quadrati CL ad OP. Ergo cuboquad.\textsuperscript{m} AD ad BI erit ut ductus cubi DC in quadratum CL ad ductum cubi IO in quadratum OP. Sed antecedentia sunt ęqualia ob latus rectum, ergo et consequentia ęqualia erunt quod promisimus primò.
\note{Comme à la vue 30, E est le grand C bouclé, distinct de C ; les couples « DE », « EI », « CE », « EO » sont lus d'après la figure de la vue 30.}

Secundum uero ostendimus sic. Recta IO ad OE est ut exponens diametralium ad differentiam exponentium, nempe ut 3 ad 2. Quare ductus cubi IO in quadratum OE erit maximus. (quod demonstrauimus iam pridem uniuersaliter geometricè sine ulla ope Algebrę) Sed ut ductus IO in OE ad ductum IC in CE ita ductus IO in OP ad IC in CL. Quare et IO in OP maximus erit.

Hęc una est ex quinque Methodis quas habemus pro tangentibus.

Lemma II

Omnia parallelogramma, puta AB, CD ad inscriptas sibi parabolas (dumodo parabolę sint eiusdem speciei) eandem rationem habent. Ponimus tantum figuram, et omittimus demonstrationem, nam facillima est, præcipuè uerò ꝑ doctrinam Indiuisibilium, addita tertia parabola in parallelogr.\textsuperscript{mo} B\uncertain{D}, et ductis ubicunq; duabus tantum rectis e f, g h. \ill{}
\note{Figure dans la marge droite : un parallélogramme vertical, A au coin supérieur droit, dont le côté gauche descend jusqu'au pied de la figure ; un parallélogramme horizontal, C au coin supérieur droit, D au pied ; ils se touchent en un point, au milieu de la figure, par où passent la droite B (horizontale) et la verticale qui descend à D. Une parabole descend du coin supérieur gauche du premier parallélogramme à ce point ; une seconde, de ce point au coin inférieur droit du second ; une troisième, en pointillé, de ce point au coin inférieur gauche. Deux droites : e f, horizontale, et g h, presque verticale. « B\uncertain{D} » : la seconde lettre est une grande lettre bouclée, qui ressemble au E de la vue 30 ; D est proposé d'après la figure, où le parallélogramme de la troisième parabole est compris entre B et D. Le dernier signe de la page, bouclé, n'est pas lu. Au pied, la réclame « Iam ».}

\page{32}
\note{Recto. En haut à droite, au-dessus de la figure, « 16 », d'un trait fin, et à sa droite un « 9 » plus gras (voir l'en-tête). Même main. La réclame « Iam » de la vue 31 est le premier mot : la lettre continue.}

Iam quadratura parabolarum talis est. Esto parabola quęlibet ABC cuius diameter BC, basis AC, parallelogrammum uerò circumscriptum sit FC. Esto Tangens AD, concipiaturq; in parallelogrammo FBDE alia parabola eiusdem speciei, siue quod idem est linea Roberualliana (noua enim illa linea quam ipse excogitasti Vir Clar.\textsuperscript{me} si ortum ducat ex aliqua parabolarum semper parabola euenit eiusdem speciei, quod ego nouum esse scio, licet fortasse turpe uideatur hoc fateri) et erit parabola ABC ęqualis trilineo ABE ut ego demonstraui ellapso iam biennio ꝑ eandem demonstrationem quam ad tuas lineas applicatam in folio separato ad te nunc mitto.
\note{Figure dans la marge droite : un rectangle, E en haut à gauche, D en haut à droite, A en bas à gauche, C en bas à droite (lettre tracée comme un G) ; F au milieu du côté gauche, B au milieu du côté droit, FB menée entre eux. Une courbe descend de E à B ; une autre, la parabole, monte de A à B ; la droite AD, la tangente, coupe le rectangle en diagonale. Comme aux vues 30–31, E est le grand C bouclé.}

Iam ut exponens applicatarum ad exponentem diametralium ita est recta DC ad CB ob tangentem ꝑ lemma p\textsuperscript{m}. Et ita est parallelogrammum EC ad FC. Sed ita etiam sunt ambę simul parabolę EBD, ABC ad parabolam ABC ob 2\textsuperscript{m} lemma. Ergo ꝑ 19 Quinti reliquum EBA siue parabola ABC ipsi ęqualis, erit ad reliquum FBA ut erat totum ad totum, hoc est ut exponens applicatar.\textsuperscript{m} ad exponentem diamet.\textsuperscript{m} Propterea conuertendo, componendoq; patet parallelogrammum quodlibet ad suam parabolam esse ut aggregatum exponentium ad exponentem applicatarum.
\note{« p\textsuperscript{m} » : « primum ». « ꝑ 19 Quinti » : la proposition 19 du livre V des Éléments, citée par l'auteur.}

Sed et alia ratione concludimus maiori breuitate quadraturam
\note{La phrase continue à la page suivante ; le bas de la page est blanc.}

\page{33}
\note{Verso de la feuille « 16 ». En haut à gauche, à l'encre grasse, « 10 » (voir l'en-tête) ; à côté transparaissent à l'envers le « 16 » et le « 9 » du recto. Même main. La phrase de la vue 32 continue ici.}

omnium parabolarum. Estq; hęc demonstratio ueluti Corollarium demonstrationis, quam in folio separato ad te mitto. Videbis uir Clar.\textsuperscript{me} ex eodem inuento ex quo gloriaris te quadraturam unius parabolę quadraticę deduxisse, nos quadraturam omnium nullo planè negotio deriuare ꝑ indiuisibilia more nostro. Quodq; elegantius est demonstratio concludit quamquam ignotum sit cuius nam speciei sit noua illa linea tua, quam diximus supra parabolam esse eiusdem gradus cum genitrice.

Esto parabolarum aliqua BAFC cuius diameter AB basis uero BC parallelogrammum circumscriptum BL, linea noua Roberualliana quęcunq; tandem sit, esto ADE, et sumpto quolibet puncto F in linea parabolica, sit applicata FI, Tangens FH, et ipsa GFD diametro ęquidistet : tum iuncta DH applicatis ęquidistans et ipsa erit ob definitionem figurę.
\note{Figure dans la marge droite : le parallélogramme BL, B en bas à gauche, C en bas à droite, L à droite en haut, A sur le côté gauche ; la parabole descend de A par F jusqu'à C ; la verticale GFD, avec O à sa rencontre de AL ; l'ordonnée FI ; la tangente FH, qui monte de F vers H, en haut à gauche ; HD en pointillé ; la ligne nouvelle, en pointillé, monte de A par D vers E, en haut à droite, hors du parallélogramme. E est de nouveau le grand C bouclé du texte ; la figure porte un e minuscule.}

Iam DF recta ad FO erit ut HI, ad IA, siue ob tangentem, ut exponens applicatarum ad exponentem diametralium. Et semper ita continget ubicumq; sumpta fuerit recta
\note{Au pied, la réclame « DF ».}

\page{34}
\note{Recto. En haut à droite, « 17 », d'un trait fin, et à sa droite « 11 », plus gras, barré d'un trait (voir l'en-tête). Même main. La réclame « DF » de la vue 33 est le premier mot : la lettre continue.}

DF. Ergo omnes antecedentes simul, nempe trilineum EDAFC siue parabola ABC ipsi ęqualis, erunt ad omnes consequentes simul, nempe ad trilineum LAC ut exponens applicatarum ad exponentem diametralium. Quadratura ergo patet conuertendo, componendoq; ut supra.

Prędictę Methodi tum pro Quadraturis, tum pro Tangentibus sunt quas minimi prę cęteris ego facio (non tamen patiar mihi illas eripi). Supersunt adhuc pro Tangentibus quatuor, pro quadraturis uero adhuc tres Methodi, inter quas una est pręcipua, et generalissima, quę communis est et parabolarum, et hyperbolarum tam quadraturis, quam etiam tangentibus, alijsq; Problematibus, quę \struck{\ill{}} nundum circumferre est animus.
\note{Le mot biffé, noyé sous l'encre, n'est pas lu ; « nundum » est écrit au-dessus.}

Iam examinauimus figuras planas, solidasq; in infinitum recta uia excurrentes. Alium modum habet Geometria in infinitum abeundi, et quidem mensurabilem. Problema circa lineam est, quę cum sit curua, et ex infinitis spiris conflata, rectę tamen lineę ęqualis est ; et hoc Geometricè more ueterum demonstratum est a nobis. Figuras quasdam huius Problematis sine ulla definitione misimus in Galliam iam ante biennium. Clar.\textsuperscript{mo} tamen Mersenno dum esset Romę communicauimus quasdam rationes pro
\note{Après « mensurabilem », un petit signe en astérisque, auquel ne répond aucune note visible sur la page. Les mots « et hoc Geometricè more ueterum demonstratum est » sont soulignés. La phrase continue à la page suivante ; le bas de la feuille est blanc.}

\page{35}
\note{Verso de la feuille « 17 », sans numéro ; même main. La phrase de la vue 34 continue ici. Cachet rond « BIBLIOTHEQUE NATIONALE » avec « R.F. » sous la date.}

descriptione proxima earumdem spiralium, quarum quędam Theoremata una cum definitione ad uos nunc mitto. Videbatur aliquid deesse, cum reperta solidorum et spatiorum in infinitum abeuntium dimensione, nulla de lineis adhuc facta esset mentio. Linearum quidem genus non potest in infinitam distantiam abire quin et quantitas infinita sit. At potest occultum haberi caput et inter laberintos infinitarum reuolutionum adeo reconditum ut solis Geometrię oculis sit uisibile. Quicquid est, et an Problema dignum sit oculis tuis ipse uidebis Vir Clarissime.

Hyperbolarum Theoremata, quę mitto ad Ill.\textsuperscript{m} De Fermat. ut iudicium subeant, num cum parabolis saltem aliqua ex parte conferri possint, uidere poteris. Si unius hyperbolę Primarię Quadratura tam diu quęsita est, Nos pro una infinitas damus. Interea toto affectu me tibi commendo V. Clarissime Vale. Dat. Florentię die \uncertain{7} Julij ann. 1646
\note{Fin de la lettre commencée à la vue 24 (« Clar.\textsuperscript{mo} Viro Roberuallio / Torr. S.P.D. ») ; aucune signature sous la date. Le chiffre du jour a la forme en chevron déjà rencontrée à la vue 26 (« 7 ad 5 ») ; il est lu 7 avec doute. Un long trait de plume sous la date, en travers du cachet.}

Factum etiam est ut usque ad initium Augusti mensis litterę iā absolutę apud me morarentur ob uaria impedimenta mea.
\note{Post-scriptum, d'une encre plus pâle et d'une écriture plus rapide et plus petite, peut-être la même main. « iā » : « iam », avec le tilde de la page.}

\page{36}
\note{Recto d'une nouvelle feuille, qui commence une nouvelle pièce. En haut à droite, « 18 », d'un trait fin (voir l'en-tête) ; pas de second nombre. Le texte est d'une main posée et régulière, à capitales de type imprimé pour les lettres des figures, que l'on rapproche de celle des vues 21–23 (rapprochement, non certitude). Cachet rond « BIBLIOTHEQUE NATIONALE » avec « R.F. ». Aucune figure sur la page. Le bas de la feuille est blanc.}

\note{En haut à gauche, d'une autre main, cursive, à l'encre noire, une note de six lignes, écrite avant ou après le titre, qu'elle serre : « Hæc est demonstratio, \struck{\ill{}} quam Torricellius in pag. 9.\textsuperscript{a} et pag. 10 epistolæ ad Roberuallium datæ 7.\textsuperscript{o} die Julii, 1646. dicit à se missam fuisse in folio separato. » Dans « epistolæ » et « datæ », la finale « æ » est ajoutée au-dessus de la ligne ou en surcharge ; le premier mot de la deuxième ligne est biffé ; après « pag. 10 » la grande capitale du titre recouvre la fin de la ligne. C'est cette note, et non le texte de la pièce, qui nomme Torricelli ; les pages 9 et 10 qu'elle cite répondent aux numéros gras « 9 » et « 10 » des vues 32 et 33, où l'auteur de la lettre annonce une démonstration « in folio separato ».}

Cl.\textsuperscript{mo} Viro Roberuallio

Lemma Pr.\textsuperscript{um}

In figura qualis a Cl. Roberual definita est, quodlibet parallelogrammum inscriptum AB minus erit quàm FA descriptum in figura \uncertain{genitrice}. Nam ductâ tangente ꝑ A, finitâq; figura, æquale erit AB ipsi AI. Ergo patet quod erat \&.
\note{L'encre n'a pas pris au début des lignes 4 à 6 du lemme : « minus erit quàm », « genitrice », « figura, æquale », « quod » sont très pâles ; « genitrice » est lu d'après les lettres « ge » et « trice », qui restent.}

Lemma secundũ

Circumscriptum vero parallelogrammũ CD majus erit quàm aliud parallelogrammum FE, quod nempe circumscriptum erit figuræ genitrici. Nam ducta tangente per O et completâ figurâ erit CD æquale ipsi FM Ergo CD maius est quàm FE Quod erat \&.

Theorema

Sumatur primò aliqua pars figuræ definitæ : puta ABCDE. Dico æqualem esse figuræ BGC. Alias major esset vel minor. Sit primò major, Seceturque BK semper bifariam donec HE minus sit excessu. Tum inscribatur in figura mixta BCE alia figura constans ex parallelogrammis æquealtis RD, LM \&c. quorum vltimum sit IO. Eritque inscripta figura ob constructionem adhuc major spatio BGC. Quod est contra Lemma primum.

Sit deinde minor. Sectoque BK semper bifariam repertum sit HE minus defectu. Tum circumscribatur figuræ BCE
\note{La phrase continue à la page suivante.}

\page{37}
\note{Verso de la feuille « 18 », sans numéro ; même main posée que la vue 36. La phrase de la vue 36 continue ici. Le texte occupe, à partir du Theorema suivant, une colonne étroite à gauche ; la moitié droite de la page est blanche, et le recto y transparaît. Aucune figure.}

alia figura constans ex parallelogrammis æquealtis, et erit circumscripta figura adhuc minor figurâ BCE. Quod esse n\textsuperscript{o} potest. Nam eadem figura circumscripta trilineo BCE major est quàm alia quædam figura circumscripta ipsi BGC ob 2.\textsuperscript{m} Lemma. Patet ergo. quod erat \&c.
\note{« n\textsuperscript{o} » : n suivi d'une petite boucle levée, abréviation de « non ».}

Abeat iam in infinitum (quamquam hactenus dicta sufficere deberent) Dico æquales esse figuras \&c. Alias altera ipsarũ major erit. Sit major APC, et erit aliqua ipsius pars puta BGC æqualis alteri figuræ. Statim absurdum patet. Si verò ponatur major ACEF erit aliqua ipsius pars, puta BCE æqualis figuræ genetrici APC. Absurdum iam patet.

Hactenus de his nouis lineis ; et ne reliquum paginæ pereat adjciamus aliquid de Spiralibus Geometricis, quæ nisi fallor et ipsæ nouæ lineæ erunt.

Esto qualibet recta linea AB secta vtcunque in C. deinde erigatur perpendicularis CD quæ sit media proportionalis inter BC, CA. Iterum recta CE media sit inter DC, CA, Secetque bifariam angulum DCA. Amplius recta CF media inter
\note{La phrase continue à la page suivante. Le bas de la page est blanc.}

\page{38}
\note{Recto. En haut à droite, « 19 », d'un trait fin (voir l'en-tête). À gauche de la vue paraissent les fins de ligne de la page de la vue 37, qui lui fait face. Même main posée ; la phrase de la vue 37 continue ici. Comme à la vue 37, le texte laisse d'abord une large marge à gauche. Aucune figure, bien que le texte renvoie à une « sequenti figurâ ».}

DC, CE secet bifariam angulum DCE : et sic fiat semper ducendo medias proportionales quæ bifariam secent angulos, habebimusque puncta A, E, F, D, B, et quotcunque alia voluerimus per quæ transibit quædam Linea curua, quam Spiralem Geometricam appellimus

Sed aliter possumus eandem definire. Si recta linea AB (in sequenti figurâ) manente eius extremo A circumagatur in plano, æquale semper velocitate : et eodem tempore aliquod punctum C ejusdem lineæ super eâdem moueatur ea lege, vt spatia æqualibus temporibus peracta inter se proportionalia sint, eandemque proportionem seruent quam habebunt distantiæ ejusdem puncti à centro A, iterum eandem spiralem Geometricam habebimus.

Peculiare hoc habet hæc linea, quod antequam perueniat ad centrum, infinitas reuolutiones circa ipsum absoluere debebit ; attamen Longitudo ejus, non solum quo ad partes, sed et vniuersa nota est. Si fuerit centrum A, vnusque ex radijs AB, ipsa veró BC tangens, et angulus BAC rectus, erit Tangens BC æqualis vniuersæ Speciali BDA.
\note{« Speciali » est bien écrit ainsi ; le sens semble demander « Spirali ». « appellimus » est l'orthographe de la page.}

Dato vero arcu quocunque BI. facillimè ipsi recta linea æqualis abscinditur, Demonstraturque.

Ostendi porterant hæc per indiuisibilia more nostro, quam Demonstrationem negligimus tamquam minus receptam
\note{« porterant » est l'orthographe de la page. La phrase continue à la page suivante ; le bas de la page est blanc.}

\page{39}
\note{Verso de la feuille « 19 », sans numéro ; le « 19 » du recto transparaît à l'envers en haut à gauche. Même main posée ; la phrase de la vue 38 continue ici. La fin de plusieurs lignes touche le bord rogné de la feuille et y perd une lettre, restituée ici entre crochets là où le sens est certain.}

ab omnibus Adhibuimus itaque duplicem positionem mor\add{e} veterum, sequuti methodum antiquorum quæ magis ab omnibus recipitur.

Spatium quoque mensurabile est quamquam ex infiniti\add{s} reuolutionibus constet, habetque notam rationem ad spatium planum finitum et certum. Extendimus deinde eandem contemplationem in alias spirales super curuis superficiebus descriptas, quas omnes rectis quibusdam Lineis demonstramus æquales.

Deinde etiam Spiralem Archimedeam in infinitas Species, propagamus, definientes, quòd non solùm spatia peract\add{a} à puncto mobili, sed etiam dignitates spatiorum proportionales sint cum angulis peractis. Ex quâ definitione totidem oriuntur spirales quot sunt parabol\add{æ}. Rationes omnium spatiorum ad circulum demonstrauit acutissimus juuenis Michael Angelus Ricci, qui infinitas quoque hyperbolas, diuersas à nostris, et ellypses contemplatus est, et n\textsuperscript{o} spernenda Theoremata inde deriuauit. Nos verò habemus et reliqua Theoremata, De rationibus abscissarum à Tangentibus cum ipsis peripharijs, præcipuè verò vnamquamque lineam spiralem vnicuique lineæ parabolicæ æqualem esse demonstramus.
\note{« n\textsuperscript{o} » : n suivi d'une petite boucle levée, abréviation de « non ». « ellypses », « peripharijs » sont l'orthographe de la page. La pièce commencée à la vue 36 s'achève ici, sans date ni signature ; le reste de la page est blanc.}

\page{40}
\note{Recto d'une nouvelle feuille. En haut à droite, « 20 », d'un trait fin (voir l'en-tête). En haut à gauche, d'une autre main, d'un trait gris pâle (crayon ou encre très pâle), soulignée et fermée d'un trait vertical : « 7. Julii 1646 ». Le texte est une copie, de la main posée des vues 36–39, de la lettre dont la vue 24 porte le début de la main rapide : même suscription, mêmes phrases, à quelques graphies près. L'encre n'a pas pris au début de plusieurs lignes ; les lettres perdues sont restituées entre crochets là où le sens est certain. Cachet rond « BIBLIOTHEQUE NATIONALE » avec « R.F. ». À gauche de la feuille, sous son bord, paraît le texte d'une autre page, qui n'appartient pas à celle-ci. Le bas de la feuille est blanc.}

Clar.\textsuperscript{mo} Viro Roberuallio

Torr. S.P.D.

De Trochoide (esto enim quantum libet Trochoides) siue Italicum siue Gallicum Problema sit, nihil mea interest. Meum certe non est quod ad inuentionem attinet. De autore ipsius quod ego \add{ac}ceperam ab amicis illud credideram scripseramq; Vos aliter vultis ! \add{p}ro me iam licet. Hoc certissimum est (quicquid tandem feratur \add{a}b alijs) celeberrimum Galileum vsq; ad supremum vitæ diem mensuram illius figuræ ignorauisse quam ex Gallia non accepit, vbi fortasse inuenta non fuerat : illud certe profitebatur ; neque video cur demonstrationem il\ill{} si a quopiam accepisset in commune non protulisset ad gloriam suam, quamquam aliena esset. Ego fateor non adeo multis ab hinc annis demonstrationes illas me reperisse, sed proprio Marte non minus, quam à quopiam alio, siue ante me, siue post, factum sit. Si verò aliqua ex meis demonstrationib\textsuperscript{s}. conuenit cum Gallicis, primùm quod ad meam internam quietem attinet, quodq; plurimi facio, ego mihi conscius sum illas omnes ex meo reperisse, \& quicunq; me nouerit idem credet : deinde quidquid alij credant, nihil me mouet. \uncertain{Eximium} illum voluptatis fructum quem percipimus vnusquisque in inuentione veritatis, et pro quo tantum speculor, nemo à me auferet De gloria quam ꝑ contentiones, et controuersias acquirere debeam minime sollicitus sum : Propterea non tantum vnam, sed \& omnes demonstrationes illas, si quis volet concedere paratus ero, \add{d}umodo ꝑ iniuriam non eripiat. Sed de Centro grauitatis Cycloidis scis profectò V. Cl.\textsuperscript{me} demonstrationem illius misisse in Galliam (atq; vtinam non misissem) precibus Cl. Mersenni integro biennio antequam illud habere diceres, vt in vltimis tandem epistolis habuisse iam diu confiteris. In illa demonstratione mea ostendebatur à me dato Centro grauitatis \& mensura alicuius plani (quod satis erat ad
\note{« il\ill{} » : la fin du mot est sous une tache d'encre ; la vue 24 porte au même endroit « illius ». « Vos aliter vultis ! » : le signe qui suit « vultis » est un point d'exclamation ou d'interrogation ; la vue 24 écrit ensuite « ꝑ me », celle-ci « pro me ». « Eximium » : la troisième voyelle se lit plutôt e. « demonstrationib\textsuperscript{s} » : b suivi d'un s levé. La phrase s'arrête au pied de la page et continue au-delà de ce lot.}

\end{document}
